% Unidad U012. Biografía estructural completa de pi.

\label{gfp:b:chap:biografia-pi-autonoma}

\section{Séparation causale entre germe, caractère et évaluation}

La construction se divise en trois strates :
\[
 \text{sélection finie}
 \longrightarrow
 \text{construction du caractère de clôture}
 \longrightarrow
 \text{évaluation archimédienne}.
\]
Le développement décimal n’intervient pas dans la première strate. La chaîne des types
est
\begin{equation}
 104\,976
 \longrightarrow468\,\mathcal U_6
 \xrightarrow{\Psi}243\,\mathcal W_6
 \lhook\joinrel\longrightarrow\mathbb F_3^6,
 \qquad|\mathbb F_3^6|=729.
 \label{gfp:b:eq:cadena-tipica-pi-autonoma}
\end{equation}
Les cardinalités comptent, respectivement, les conditions initiales
enrichies, les émissions distinctes, les mots ternaires visibles et les éléments
de la fibre ambiante complète.

Pour \(U_6=(u_1,\ldots,u_6)\), la projection est
\[
 \Psi(U_6)=((-u_1)\bmod3,\ldots,(-u_6)\bmod3).
\]
Le recensement des \(43\) orbites de \(D_3\) démontre que l’orbite de
clôture est la seule dont les six orientations possèdent à la fois :
\[
 \text{cinq relèvements},\qquad
 \text{masse }1008,\qquad
 1008=432+4\cdot144.
\]
Le calibre
\[
 \operatorname{diag}_c=6,
 \qquad
 \operatorname{card}=ES
\]
sélectionne un unique représentant :
\begin{equation}
 \boxed{w_6^\pi=010211.}
 \label{gfp:b:eq:semilla-pi-autonoma}
\end{equation}
La sélection utilise le catalogue, les multiplicités, l’action orbitale et l’orientation ;
elle ne consulte pas un préfixe de \(\pi\).

\section{Relèvements \texorpdfstring{\(L_0,L_1\)}{L0, L1} et \(1^{\mathrm a}\) frontière coinductive}

Sur les vecteurs-lignes de \(\mathbb F_3^6\), soient
\[
L_0=
\begin{pmatrix}
2&2&2&1&2&1\\
2&1&2&2&1&1\\
1&0&0&1&0&2\\
0&0&1&1&1&0\\
2&2&2&0&2&1\\
2&1&2&1&0&0
\end{pmatrix},
\quad
L_1=
\begin{pmatrix}
2&1&2&0&2&2\\
1&2&1&1&2&0\\
1&2&1&1&1&2\\
0&1&0&0&2&1\\
2&0&2&1&0&1\\
0&2&2&2&1&1
\end{pmatrix}.
\]
Si \(b_1=w_6^\pi\), la récurrence des blocs est
\[
 b_{k+1}=b_kL_{\varepsilon_k},
 \qquad
 (\varepsilon_1,\varepsilon_2,\varepsilon_3,\varepsilon_4)=(0,0,1,1),
\]
et le mot accumulé est formé par concaténation :
\[
 w_{6(k+1)}=w_{6k}\mathbin{\|}b_{k+1}.
\]
On ne multiplie pas un mot de longueur \(6k\) par une matrice \(6\times6\).

Pour le canal de clôture :
\[
\begin{array}{c|c}
\text{profondeur}&\text{nouveau bloc}\\ \hline
6&010211\\
12&012222\\
18&010211\\
24&002111\\
30&110221
\end{array}
\]
et
\begin{equation}
 \boxed{
 w_{30}^{\pi}
 =010211|012222|010211|002111|110221.}
 \label{gfp:b:eq:w30-pi-autonomo}
\end{equation}
La frontière suivante est
\[
 u_5^\pi=222220.
\]
Avec les autres canaux,
\[
 R_{36}=
 \begin{pmatrix}
 222220\\021222\\102011
 \end{pmatrix}.
\]
Le symbole \(R_{36}\) enregistre le sixième bloc de trois histoires et ouvre une
frontière ; ce n’est pas un terminal.

La branche certifiée de vingt blocs commence par
\begin{align*}
\pi:\;&010211|012222|010211|002111|110221|222220|111201|\\
&212121|200121|100100|101222|022212|012012|111210|\\
&121011|200220|120210|000101|022010|020111.
\end{align*}
C’est pourquoi ni \(w_{30}\) ni \(R_{36}\) ne terminent la généalogie.

\section{Microfibre de clôture de \(\pi\) : \(5\) régions orientées et multiplicité \(432+4\cdot144=1008\)}

La microfibre
\[
 \mathscr R_\pi^{\rm micro}
 :=\{U\in\mathcal U_6:\Psi(U)=010211\}
\]
contient exactement cinq régions :
\[
\begin{array}{c|c|c|c}
U_6&E_{\rm sig}&\operatorname{mult}&\text{rôle}\\ \hline
501|614|498|169|272|272&555555&432&\perp\\
810|923|870|169|272|272&339555&144&++\\
810|983|810|169|272|272&393555&144&++\\
870|923|810|169|272|272&933555&144&++\\
870|923|810|418|674|521&933717&144&--
\end{array}
\]
En conséquence,
\begin{equation}
 \boxed{
 5=1_\perp+3_{++}+1_{--},
 \qquad
 1008=432+4\cdot144.}
 \label{gfp:b:eq:descomposiciones-cinco-regiones-pi}
\end{equation}
La première identité classe des fonctions d’orientation. La seconde compte
des multiplicités de provenance. Elles ne comptent pas cinq copies indiscernables.

La région transversale est
\[
 501|614|498|169|272|272,
 \qquad E_{\rm sig}=555555,
 \qquad\operatorname{mult}=432.
\]
Le retour muté est
\[
 870|923|810|418|674|521,
 \qquad E_{\rm sig}=933717,
 \qquad\operatorname{mult}=144.
\]
Les intervertir conserverait le recensement \(3/1/1\), mais fausserait la structure
régionale complète.

\section{Transfert de phase sur le multigraphe des blocs complets}

Pour \(U_6=(U_1,\ldots,U_6)\), soit
\[
 A(U_6)=(U_1,U_2,U_3),
 \qquad
 B(U_6)=(U_4,U_5,U_6).
\]
Chaque émission définit une arête entre \(A(U_6)\) et \(B(U_6)\). Le multigraphe
brut possède \(96\) sommets, répartis en composantes de tailles
\[
 28,28,28,4,4,4.
\]
On passe au quotient de la phase interne par
\[
 (a,b,c)\sim(b,c,a)\sim(c,a,b).
\]
Le représentant est la rotation lexicographiquement minimale. Le graphe
quotient possède \(32\) sommets et deux composantes de tailles \(28\) et \(4\).

Les ancrages sont
\[
\begin{aligned}
U_{\rm APP}&=903|850|850|252|109|252,\\
U_e&=850|963|890|149|252|212,\\
U_\varphi&=830|943|830|109|252|252.
\end{aligned}
\]
\Needspace{9\baselineskip}
Nous numérotons les régions
\[
\begin{aligned}
U_\pi^{(1)}&=810|923|870|169|272|272,\\
U_\pi^{(2)}&=810|983|810|169|272|272,\\
U_\pi^{(3)}&=870|923|810|169|272|272,\\
U_\pi^{(4)}&=870|923|810|418|674|521,\\
U_\pi^{(5)}&=501|614|498|169|272|272.
\end{aligned}
\]

\section{\(5\) cycles orientés de clôture de la microfibre}

\subsection{Clôture de \texorpdfstring{\(U_\pi^{(1)}\)}{U-pi 1}}
\begin{align*}
&252|109|252|903|850|850\to
903|850|850|252|109|252\to\\
&252|109|252|923|870|810\to
810|923|870|169|272|272\to\\
&272|169|272|963|890|850\to
850|963|890|149|252|212\to\\
&212|149|252|943|830|830\to
830|943|830|109|252|252\to
252|109|252|903|850|850.
\end{align*}

\Needspace{10\baselineskip}
\subsection{Clôture de \texorpdfstring{\(U_\pi^{(2)}\)}{U-pi 2}}
\begin{align*}
&903|850|850|252|109|252\to
272|169|272|903|850|850\to\\
&272|169|272|983|810|810\to
810|983|810|169|272|272\to\\
&272|169|272|963|890|850\to
850|963|890|149|252|212\to\\
&212|149|252|943|830|830\to
830|943|830|109|252|252\to
903|850|850|252|109|252.
\end{align*}

\subsection{Clôture de \texorpdfstring{\(U_\pi^{(3)}\)}{U-pi 3}}
\begin{align*}
&903|850|850|252|109|252\to
292|189|232|903|850|850\to\\
&292|189|232|923|810|870\to
870|923|810|169|272|272\to\\
&272|169|272|963|890|850\to
850|963|890|149|252|212\to\\
&212|149|252|943|830|830\to
830|943|830|109|252|252\to
903|850|850|252|109|252.
\end{align*}

\subsection{Clôture de \texorpdfstring{\(U_\pi^{(4)}\)}{U-pi 4}}
\begin{align*}
&903|850|850|252|109|252\to
272|169|272|903|850|850\to\\
&272|169|272|963|890|850\to
850|963|890|149|252|212\to\\
&212|149|252|923|810|870\to
870|923|810|418|674|521\to\\
&943|830|830|521|418|674\to
830|943|830|109|252|252\to
903|850|850|252|109|252.
\end{align*}

\subsection{Clôture de \texorpdfstring{\(U_\pi^{(5)}\)}{U-pi 5}}
\begin{align*}
&252|109|252|903|850|850\to
903|850|850|252|109|252\to\\
&614|498|501|252|109|252\to
501|614|498|169|272|272\to\\
&272|169|272|963|890|850\to
850|963|890|149|252|212\to\\
&212|149|252|943|830|830\to
830|943|830|109|252|252\to
252|109|252|903|850|850.
\end{align*}

\begin{theorem}[Minimalité des cinq clôtures]
\label{gfp:b:thm:minimalidad-cierres-pi-autonoma}
Pour chaque \(i\in\{1,\ldots,5\}\), la longueur minimale d’une marche fermée
qui parcourt simultanément
\[
 U_\pi^{(i)},\quad U_e,\quad U_\varphi,\quad U_{\rm APP}
\]
dans le multigraphe quotient est huit.
\end{theorem}

\begin{proof}
Pour \(i\) fixé, soit
\[
 E_*=\{U_\pi^{(i)},U_e,U_\varphi,U_{\rm APP}\}.
\]
On considère l’espace fini
\[
 \mathscr S_i=\{(v,M):v\in V(G),\ M\subseteq E_*\}.
\]
La première composante enregistre le sommet et la seconde les arêtes cibles
déjà parcourues. Si une arête \(e\) mène de \(v\) à \(v'\), la transition est
\[
 (v,M)\longmapsto
 \bigl(v',M\cup(E_*\cap\{e\})\bigr).
\]
La recherche en largeur énumère les longueurs par ordre croissant. Le premier
retour à \((v,E_*)\) apparaît à la longueur huit pour les cinq régions.
Les marches imprimées réalisent la borne ; l’optimalité de la recherche exclut
une longueur inférieure.
\end{proof}

La fonction de clôture est encodée par
\[
 \mathfrak C_\pi=
 \bigl(\mathscr R_\pi^{\rm micro},
 w_\pi,\rho_\pi,m_\pi,\ell_{\rm cl}\bigr),
\]
\[
 \rho_\pi=(\perp,++ ,++ ,++ ,--),
 \qquad
 m_\pi=(432,144,144,144,144),
 \qquad
 \ell_{\rm cl}=8.
\]

\section{Opérateur d’incidence et demi-tour orienté}

L’ombre d’incidence est
\[
 A_W=
 \begin{pmatrix}
 0&1&1&1&1&1\\
 1&0&1&2&2&1\\
 1&1&0&1&2&2\\
 1&2&1&0&1&2\\
 1&2&2&1&0&1\\
 1&1&2&2&1&0
 \end{pmatrix}.
\]
La multiplication dans \(\mathbb F_3\) donne
\begin{equation}
 \boxed{A_W^2=2I_6=-I_6.}
 \label{gfp:b:eq:aw-cuadrado-menos-identidad-pi}
\end{equation}
Les produits entre lignes et colonnes distinctes s’annulent modulo trois,
tandis que chaque produit diagonal vaut deux. Une application réalise un quart
de tour ternaire ; deux applications réalisent l’inversion orientée.

\section{Pente \texorpdfstring{\(1/5\)}{1/5} déterminée par la partition régionale \(432+4\cdot144=1008\)}

Soit
\[
 \lambda=(\lambda_1,\ldots,\lambda_5)^{\mathsf T}
 \in\mathbb Q_{\geq0}^5,
 \qquad
 \mathbf1^{\mathsf T}\lambda=1.
\]
Oublier l’étiquette régionale applique le symétriseur
\[
 P_5=\frac15\mathbf1\mathbf1^{\mathsf T}.
\]
Pour tout vecteur normalisé,
\[
 P_5\lambda=\frac15\mathbf1.
\]
L’invariance \(P_5\lambda=\lambda\) a donc une unique solution :
\begin{equation}
 \boxed{\lambda=\frac15\mathbf1,\qquad t_1=\frac15.}
 \label{gfp:b:eq:pendiente-primitiva-pi}
\end{equation}
L’entier cinq provient de la microfibre régionale ; il n’est pas reconnu dans une
formule de \(\pi\).

\section{Composition rationnelle et compensateur \texorpdfstring{\(239\)}{239}}

La loi de composition des pentes est
\[
 t\star s:=\frac{t+s}{1-ts}.
\]
Le premier doublement donne
\[
 t_2=t_1\star t_1
 =\frac{2/5}{1-1/25}
 =\frac5{12}.
\]
Le second donne
\[
 t_4=t_2\star t_2
 =\frac{10/12}{1-25/144}
 =\frac{120}{119}.
\]
Le retour à la diagonale unitaire exige \(q>0\) tel que
\[
 \frac{120/119-1/q}{1+(120/119)/q}=1.
\]
La multiplication par \(119q\) produit
\[
 120q-119=119q+120,
\]
et impose
\begin{equation}
 \boxed{q=239.}
 \label{gfp:b:eq:compensador-239-pi}
\end{equation}

\section{Identité gaussienne et \(1^{\mathrm a}\) demi-tour orienté}

On définit
\begin{equation}
 \Theta_\pi
 :=16\arctan\frac15-4\arctan\frac1{239}.
 \label{gfp:b:eq:theta-pi-autonoma}
\end{equation}
Les carrés successifs sont
\[
\begin{array}{c|r@{}l}
\text{puissance}&\multicolumn{2}{c}{\text{entier gaussien}}\\ \hline
(5+i)^2&24&+10i\\
(5+i)^4&476&+480i\\
(5+i)^8&-3824&+456960i\\
(5+i)^{16}&-208797818624&-3494830080i\\
(239-i)^2&57120&-478i\\
(239-i)^4&3262465916&-54606720i
\end{array}
\]
et la multiplication finale donne
\begin{equation}
 \boxed{
 (5+i)^{16}(239-i)^4
 =-681386607803576157184.}
 \label{gfp:b:eq:identidad-gaussiana-pi-autonoma}
\end{equation}
La partie imaginaire est nulle et la partie réelle est négative. Par conséquent,
\[
 \Theta_\pi\equiv\pi\pmod{2\pi}.
\]
Para \(0<x<1\),
\[
 x-\frac{x^3}{3}<\arctan x<x.
\]
D’où
\[
 \Theta_\pi>
 16\left(\frac15-\frac1{3\cdot5^3}\right)-\frac4{239}>3
\]
et
\[
 \Theta_\pi<\frac{16}{5}<4.
\]
L’unique orientation négative dans cet intervalle est le premier demi-tour
positif :
\begin{equation}
 \boxed{\Theta_\pi=\pi.}
 \label{gfp:b:eq:reconocimiento-pi-exacto}
\end{equation}

L’ordre de la reconnaissance archimédienne ultérieure est
\[
 \text{cinq régions}
 \to\frac15
 \to\frac5{12}
 \to\frac{120}{119}
 \to239
 \to\Theta_\pi
 \to\pi.
\]

\section{Encadrements rationnels de précision arbitraire}

Pour \(0<x<1\), soit
\[
 S_n(x)=\sum_{r=0}^{n}(-1)^r\frac{x^{2r+1}}{2r+1}.
\]
Le critère de Leibniz donne
\[
 S_{2m+1}(x)<\arctan x<S_{2m}(x).
\]
On définit
\[
 a_m^-=S_{2m+1}(1/5),\quad a_m^+=S_{2m}(1/5),
\]
\[
 b_m^-=S_{2m+1}(1/239),\quad b_m^+=S_{2m}(1/239),
\]
et
\begin{equation}
 \ell_{\pi,m}=16a_m^- -4b_m^+,
 \qquad
 h_{\pi,m}=16a_m^+ -4b_m^-.
 \label{gfp:b:eq:horquilla-pi-autonoma}
\end{equation}
Alors
\[
 \ell_{\pi,m}<\pi<h_{\pi,m}
\]
et
\begin{equation}
 h_{\pi,m}-\ell_{\pi,m}
 =
 \frac{16}{(4m+3)5^{4m+3}}
 +\frac{4}{(4m+3)239^{4m+3}}
 \longrightarrow0.
 \label{gfp:b:eq:anchura-horquilla-pi}
\end{equation}
Les bornes inférieures croissent et les bornes supérieures décroissent.

Comme \(3<\pi<4\), le lecteur ternaire agit sur
\[
 r_\pi:=\pi-3\in(0,1),
\]
avec encadrement
\[
 [\ell_{\pi,m}-3,h_{\pi,m}-3].
\]

\section{Certification archimédienne ultérieure de la fibre déjà engendrée}

Supposons que la relation interne \(\operatorname{Ext}\) et la survie
coinductive aient déjà produit, sans utiliser Machin ni un développement tabulé,
l’extension compatible \(N_{K+1}\) du préfixe \(N_K\), de longueur
\(L_K=6K\). Les \(729\) sous-cylindres de sa fibre ambiante sont
\[
 I_{K,u}=
 \left[
 \frac{729N_K+u}{3^{L_K+6}},
 \frac{729N_K+u+1}{3^{L_K+6}}
 \right),
 \qquad0\leq u<729.
\]
L’ordre \(m=m(K)\) est ensuite augmenté jusqu’à ce que l’encadrement de Leibniz
localise le sous-cylindre déjà produit. Soient
\[
 \ell_{\pi,K}:=\ell_{\pi,m(K)}-3,
 \qquad
 h_{\pi,K}:=h_{\pi,m(K)}-3.
\]
Les extrémités entières compatibles sont
\[
 j_{\min}=
 \max\!\left(
 729N_K,\left\lfloor
 \ell_{\pi,K}3^{L_K+6}\right\rfloor
 \right),
\]
\[
 j_{\max}=
 \min\!\left(
 729N_K+728,\left\lceil
 h_{\pi,K}3^{L_K+6}\right\rceil-1
 \right).
\]
La condition de localisation unitaire est
\begin{equation}
 \boxed{j_{\min}=j_{\max},\qquad N_{K+1}=j_{\min}.}
 \label{gfp:b:eq:seleccion-unitaria-subcilindro-pi}
\end{equation}
La première égalité certifie la localisation unitaire ; la seconde vérifie
que cette localisation coïncide avec l’extension engendrée intérieurement. Aucune
des deux ne définit la transition TPK ni ne sélectionne parmi les \(729\) états.
L’existence et l’unicité procèdent de \(\operatorname{Ext}\) et de la
survie ; Machin et les encadrements de Leibniz sont des reconnaissances
ultérieures.

\section{Prolongement coinductif de la section de \(\pi\) par le système inverse des histoires survivantes}

Soit \(\mathcal H_m^\pi\) l’ensemble des histoires enrichies à la profondeur
\(m\), et soit
\[
 p_{n,m}:\mathcal H_n^\pi\longrightarrow\mathcal H_m^\pi
\]
la troncature. Nous définissons
\[
 \operatorname{Surv}_m^{(N),\pi}
 :=p_{N,m}(\mathcal H_N^\pi),
\]
\[
 \operatorname{Surv}_m^\pi
 :=\bigcap_{N\geq m}\operatorname{Surv}_m^{(N),\pi}.
\]
Lorsque la relation \(\operatorname{Ext}\) est non vide, univoque et
compatible à tous les niveaux, les ensembles sont des singletons et satisfont
\[
 p_{m+1,m}(\operatorname{Surv}_{m+1}^{\pi})
 =\operatorname{Surv}_{m}^{\pi}.
\]
Dans ce domaine de prolongeabilité, il existe donc une unique section
\[
 x_\pi\in
 X_\pi:=\varprojlim_m\operatorname{Surv}_m^\pi.
\]
Les cylindres sont emboîtés et leurs diamètres sont \(3^{-6K}\), qui tend vers
zéro. Leur intersection contient un unique point, identifié par le caractère
angulaire à \(r_\pi=\pi-3\).

Pour toute profondeur décimale \(M\), on peut choisir un niveau \(K(M)\) dont le
cylindre est contenu dans une seule cellule décimale de diamètre \(10^{-M}\).
Les préfixes profonds se tronquent exactement aux précédents.

\section{Caractérisation structurelle, conditions nécessaires de validité et obstructions}

\begin{theorem}[Caractérisation du canal de clôture]
\label{gfp:b:thm:caracterizacion-final-pi-autonoma}
À partir du catalogue, de l’action \(D_3\), de la jauge orientée, des
opérateurs \(L_0,L_1,A_W\), de la loi des pentes et de la filtration rationnelle :
\begin{enumerate}
 \item on sélectionne \(w_6^\pi=010211\) ;
 \item on construit cinq régions avec les identités de
       \eqref{gfp:b:eq:descomposiciones-cinco-regiones-pi} ;
 \item chaque région appartient à une clôture minimale de longueur huit ;
 \item \(A_W^2=-I_6\) fixe l’inversion orientée ;
 \item la symétrisation régionale force \(1/5\) ;
 \item la composition force \(120/119\) ;
 \item le retour unitaire force \(239\) ;
 \item l’identité gaussienne reconnaît le premier demi-tour ;
 \item les encadrements et la fibre de \(729\) éléments produisent une section
       compatible à toute profondeur finie.
\end{enumerate}
Son évaluation est
\[
 \boxed{\pi=16\arctan\frac15-4\arctan\frac1{239}.}
\]
\end{theorem}

La chaîne complète est
\[
 \mathrm{APP}\to\mathrm{TRIT}\to\mathrm{TPK}
 \to104\,976\to468\to243\to w_6^\pi
 \to w_{30}^\pi\to R_{36}\to\Gamma_9
 \to X_\pi\to\pi.
\]

\begin{remark}[Certificat et contrôle de validité : réfutation du canal de clôture]
La dérivation est réfutée si :
\begin{enumerate}
 \item l’orbite de clôture cesse d’être unique ;
 \item l’une des cinq régions disparaît ou son identité complète change ;
 \item l’une des cinq clôtures a une longueur inférieure à huit ;
 \item \(A_W^2\ne2I_6\);
 \item le symétriseur admet une autre distribution normalisée ;
 \item les compositions ne produisent pas \(5/12\), \(120/119\) et \(239\) ;
 \item l’identité gaussienne échoue ;
 \item un encadrement cesse de contenir le caractère angulaire ;
 \item un niveau ne parcourt pas les \(729\) éléments de la fibre ambiante ou ne
       localise pas exactement une extension ;
 \item un développement cible intervient avant la reconnaissance.
\end{enumerate}
\end{remark}
